\QHeading{第1题　下面是某地区一周的日平均气温统计表。}
\par\noindent


\DataTable{2}{星期 & 平均气温/$^\circ$C\\\hline
一 & 15\\\hline
二 & 19\\\hline
三 & 21\\\hline
四 & 23\\\hline
五 & 18\\\hline
六 & 17\\\hline
日 & 20\\\hline}
\par

\PlainItem{① 这周的日平均气温最高是（\textcolor{answerred}{\ensuremath{23}}）℃，最低是（\textcolor{answerred}{\ensuremath{15}}）℃。}
\Problem{② 请你算一算，这周的平均气温是多少摄氏度？}

\par\noindent

\QHeading{第1题②　补充方法：用19℃作基数}

\Answer{高于19℃的部分与低于19℃的部分正好相抵。多出量为2、4、1；不足量为4、1、2。}

\QHeading{第2题　下面是乐乐家和明明家1～3月份的用电量统计表。}

\Problem{乐乐家1～3月份的用电量统计表}
\par

\DataTable{2}{月份 & 用电量/\newline 千瓦时\\\hline
1 & 58\\\hline
2 & 86\\\hline
3 & 63\\\hline
平均每月\newline 用电量 & {\color{answerred}$69$}\\\hline}
\Problem{明明家1～3月份的用电量统计表}
\DataTable{2}{月份 & 用电量/\newline 千瓦时\\\hline
1 & 62\\\hline
2 & 92\\\hline
3 & 65\\\hline
平均每月\newline 用电量 & {\color{answerred}$73$}\\\hline}
\Problem{① 求出平均每月用电量并填在表中。② 将上面的两个表合并成一个表。乐乐、明明两家1～3月份的用电量统计表}
\DataTable{3}{月份 & 乐乐家用电量/千瓦时 & 明明家用电量/千瓦时\\\hline
1 & {\color{answerred}$58$} & {\color{answerred}$62$}\\\hline
2 & {\color{answerred}$86$} & {\color{answerred}$92$}\\\hline
3 & {\color{answerred}$63$} & {\color{answerred}$65$}\\\hline
平均每月用电量 & {\color{answerred}$69$} & {\color{answerred}$73$}\\\hline}
\Answer{答：乐乐家平均每月用电69千瓦时。}
\Answer{答：明明家平均每月用电73千瓦时。}
\par\noindent

\QHeading{第2题①　补充方法：乐乐家用电量}

\Answer{用60千瓦时作基数，先平衡三个月的多用量与少用量。多出量为26、3；不足量为2。}

\QHeading{第2题①　补充方法：明明家用电量}

\Answer{用70千瓦时作基数，先平衡三个月的多用量与少用量。多出量为22；不足量为8、5。}

\QHeading{第3题　解决问题。}

\Problem{小伟在期末测试中语文、数学、英语三科的平均分是96分，其中语文得了90分，英语得了98分。小伟在期末测试中数学得了多少分？}
\par

\Answer{{\Formula{\begin{aligned}[t]&\phantom{={}}96\times3-90-98\\&=288-90-98\\&=198-98\\&=100\text{（分）}\end{aligned}}}}
\Answer{答：数学考了100分。}



\QHeading{多出量合计}

\Answer{
\Formula{\begin{aligned}[t]&\phantom{={}}2+4+1\\
&=6+1\\
&=7\text{（℃）}\end{aligned}}}

\QHeading{不足量合计}

\Answer{
\Formula{\begin{aligned}[t]&\phantom{={}}4+1+2\\
&=5+2\\
&=7\text{（℃）}\end{aligned}}}

\QHeading{计算平均数}

\Answer{
\Formula{\begin{aligned}[t]&\phantom{={}}19+(7-7)\div 7\\
&=19+0\div 7\\
&=19+0\\
&=19\text{（℃）}\end{aligned}}\par 答：这周平均气温是19℃。}

\QHeading{多出量合计}

\Answer{
\Formula{\begin{aligned}[t]&\phantom{={}}26+3\\
&=29\text{（千瓦时）}\end{aligned}}}

\QHeading{计算平均数}

\Answer{
\Formula{\begin{aligned}[t]&\phantom{={}}60+(29-2)\div 3\\
&=60+27\div 3\\
&=60+9\\
&=69\text{（千瓦时）}\end{aligned}}\par 答：乐乐家平均每月用电69千瓦时。}

\QHeading{不足量合计}

\Answer{
\Formula{\begin{aligned}[t]&\phantom{={}}8+5\\
&=13\text{（千瓦时）}\end{aligned}}}

\QHeading{计算平均数}

\Answer{
\Formula{\begin{aligned}[t]&\phantom{={}}70+(22-13)\div 3\\
&=70+9\div 3\\
&=70+3\\
&=73\text{（千瓦时）}\end{aligned}}\par 答：明明家平均每月用电73千瓦时。}

\QHeading{第3题　补充方法：用平均分补差}

\Answer{语文比96分少6分，英语比96分多2分；数学必须再多4分才能平衡。
\Formula{\begin{aligned}[t]&\phantom{={}}96+[(96-90)-(98-96)]\\
&=96+[6-2]\\
&=96+4\\
&=100\text{（分）}\end{aligned}}\par 答：小伟数学得了100分。}

\Answer{
\QHeading{计算过程：第1题：一周的平均气温}
\Formula{
\begin{aligned}[t]&\phantom{={}}(15+19+21+23+18+17+20)\div7\\&=(34+21+23+18+17+20)\div7\\&=(55+23+18+17+20)\div7\\&=(78+18+17+20)\div7\\&=(96+17+20)\div7\\&=(113+20)\div7\\&=133\div7\\&=19\text{（℃）}\end{aligned}
}
答：这周平均气温为19℃。

\QHeading{计算过程：第2题：两家平均每月用电量}
乐乐家：
\Formula{
\begin{aligned}[t]&\phantom{={}}(58+86+63)\div3\\&=(144+63)\div3\\&=207\div3\\&=69\text{（千瓦时）}\end{aligned}
}
答：乐乐家平均每月用电69千瓦时。

明明家：
\Formula{
\begin{aligned}[t]&\phantom{={}}(62+92+65)\div3\\&=(154+65)\div3\\&=219\div3\\&=73\text{（千瓦时）}\end{aligned}
}
答：明明家平均每月用电73千瓦时。

}
